\begin{textsample}{POS Dim 3 – am – Score 29.00 – am/am\_inf\_10.txt}  \label{ex:f3_pos_032}
Interações e \textbf{brincadeira} As interações e as \textbf{brincadeiras} dos profissionais com os \textbf{bebês} e \textbf{crianças} devem acontecer a todo o momento no processo de aprendizagem e desenvolvimento, favorecendo o diálogo e estabelecendo relações ricas e desafiadoras com o mundo da cultura.
As múltiplas linguagens vividas nas \textbf{creches} e \textbf{pré-escolas} por meio das interações com \textbf{crianças} e \textbf{adultos} contribuem para a construção da própria linguagem, ajudando -as a expressarem seus \textbf{desejos} e anseios, assim como as hipóteses que elas elaboram sobre o mundo.
É indispensável que a/o professora/professor da \textbf{creche} interaja com as \textbf{crianças} com a intenção de entender e responder as tentativas de comunicação por meio de \textbf{sons}, \textbf{gestos} e palavras, valorizando o diálogo que se intensificará ao longo do tempo, por isso a necessidade da/do professora/professor estar atento a cada ação da \textbf{criança}.
Na Pré-Escola, a ampliação da linguagem deve ser proporcionada pelos espaços onde a fala da \textbf{criança} assuma centralidade e importância.
Estimular a \textbf{criança} a se comunicar deve ser uma preocupação constante da/do professora/professor, que precisa garantir tempo para que a \textbf{criança} se expresse.
Materiais pedagógicos como fantoches, telefones, microfones, livros de literatura, revistas, cartões com imagens, CD e DVD com material de boa qualidade, entre outros, ampliam o repertório dos \textbf{bebês} e \textbf{crianças}, criando necessidade de expressar as apreensões que fazem daquilo que lhes \textbf{interessam}.
A literatura é de grande valor nesse processo.
É necessário que haja uma variedade de livros em bom estado e adequados à faixa etária, considerando a qualidade literária dos textos escritos e imagens.
A \textbf{criança} precisa ter liberdade para explorar os livros e estes precisam estar acessíveis, devendo ser utilizados cotidianamente nas atividades, possibilitando um amplo desenvolvimento linguístico a partir de histórias narradas pelo professor, assim como as relatadas pelas \textbf{crianças} durante os momentos de interação com a turma.
Cabe a/ao professora/professor promover momentos agradáveis que proporcionem ao grupo interações afetuosas, empáticas e respeitosas.
A atividade que mais desenvolve o \textbf{bebê} é a comunicação emocional ( ELKONIN, 2017 ; LÍSINA, 1987 ). \textbf{Conversar} com o \textbf{bebê} com atenção, destinando tempo e \textbf{cuidado} para esse momento é essencial para o desenvolvimento das funções psíquicas em formação.
Assim, não se deve ignorar uma comunicação de qualidade com o \textbf{bebê} porque ele não fala.
Deve -se \textbf{conversar} sobre os \textbf{cuidados} que naquele momento o \textbf{bebê} vai receber, \textbf{contar} histórias, ler, cantar, dançar e explicar as propostas.
É nesse processo de interação com pessoas falantes que o \textbf{bebê} vai estabelecer formas de compreensão de funcionamento da linguagem, caminhando \textbf{progressivamente} para o uso de formas de comunicação próprias do seu desenvolvimento, como o \textbf{gesto}.
A linguagem deve proporcionar à \textbf{criança} o desenvolvimento do pensamento.
A \textbf{criança} deve ser questionada sempre que necessário sobre o seu \textbf{brincar}, sobre o porquê das coisas que faz e sobre a \textbf{rotina} da sala de referência[^1 ], assim como sobre seus sentimentos.
A \textbf{criança} precisa da \textbf{ajuda} dos \textbf{adultos} para identificar os sentimentos e como lidar com eles.
Neste contexto, a linguagem informal deve ser respeitada.
O acesso à linguagem formal se dá na convivência com pessoas falantes e ela não deve ser ensinada de forma agressiva ou humilhante.
O maior desafio é que a \textbf{criança} se expresse, fazendo uso das variadas linguagens.
O \textbf{brincar} se constitui em eixo norteador de todo o processo da Educação \textbf{Infantil}.
Sua importância deve ser visualizada no tempo e espaço destinados a essa atividade.
Livre ou orientada, no coletivo, individual e/ou com a participação dos profissionais, o \textbf{brincar} é elemento que constitui a cultura \textbf{infantil}.
Por \textbf{volta} dos dois anos, os objetos passam a ser o centro da atenção das \textbf{crianças} bem \textbf{pequenas}.
É importante oferecer \textbf{brinquedos} estruturados, não estruturados, móbiles, \textbf{brinquedos} \textbf{sensoriais} feitos com materiais reaproveitados, como panelas para \textbf{bater} e garrafas para colocar coisas e líquidos diferentes.
Atentando sempre para a segurança, estes objetos fazem parte da atividade que mais desenvolve a \textbf{criança} bem \textbf{pequena} : a atividade de manipulação objetal ( ELKONIN, 2017 ).
Durante as \textbf{brincadeiras} coletivas, orienta -se a atenção e/ou mediação por parte do professor no sentido de preservar a integralidade física.
É ainda necessário ressaltar a importância da compreensão das funções das \textbf{brincadeiras} e dos \textbf{brinquedos}, pois com eles, a \textbf{criança} se relaciona com seus pares, exercita a linguagem, o pensamento e a memória, amplia o movimento, comunica sentimentos, imaginação e lida com diferentes emoções ( alegria, satisfação, frustração, raiva ). \textbf{Brincar} é uma atividade cultural, portanto, aprende -se a \textbf{brincar} com os pares, com parceiros mais experientes e com os \textbf{adultos}.
Na Pré-Escola, a \textbf{brincadeira} de faz de conta se torna a atividade principal do desenvolvimento psíquico.
Nessa faixa etária, a \textbf{criança} se \textbf{interessa} pelas relações sociais dos \textbf{adultos} e pela \textbf{brincadeira}, pode assumir diferentes papéis, inicialmente \textbf{imitando} para depois, representar.
Durante as \textbf{brincadeiras} de faz de conta, a \textbf{criança} passa a aprender e experimentar diferentes papéis e ações, e isso favorece o entendimento das regras sociais, o desenvolvimento da memória, da percepção e do controle da conduta.
O \textbf{brincar} de faz de conta proporciona o desenvolvimento da linguagem, incentivando também a comunicação com ampliação do vocabulário, da expressão e do exercício da compreensão do outro.
As \textbf{brincadeiras} externas com \textbf{areia} e água são divertidas e prazerosas, pois é por meio das descobertas que a \textbf{criança} vai ampliando sua experiência com o mundo.
O tanque de \textbf{areia} é um exemplo de instrumento que estimula os sentidos da \textbf{criança}.
Este pode ser enriquecido com \textbf{brinquedos}, favorecendo o uso da \textbf{areia} de diferentes formas e maneiras, como experimentá -la seca, úmida e molhada.
Para que este tanque de \textbf{areia} continue sendo para a \textbf{criança} um recurso importante, são necessários alguns \textbf{cuidados} de \textbf{limpeza} \textbf{diária}, como cobri -lo com lona durante a noite, evitando chuva e animais e observar a existência de pedras em meio à \textbf{areia}, evitando que a \textbf{criança} possa ingerir ou se machucar.
Para que seja utilizado com eficácia, a \textbf{areia} deve ser previamente tratada.
A/o professora/professor de Educação \textbf{Infantil} deve \textbf{brincar} com as \textbf{crianças}, \textbf{conversar}, ouvir e \textbf{contar} histórias, rabiscar, desenhar, modelar, construir, pintar, sentar no \textbf{chão}, comer, rir, auxiliar nas atividades, propor situações desafiadoras, participar e interagir na prática com a \textbf{criança}, \textbf{demonstrando} segurança no seu fazer pedagógico.
O compromisso da/do professora/professor com a garantia dos direitos de aprendizagem e desenvolvimento, efetivados por meio dos objetivos dos campos de experiências deve ser objetivo primordial do seu planejamento. \textbf{Brincar} é a principal linguagem da infância, por isso, as \textbf{crianças} precisam \textbf{brincar} independentemente de suas condições motoras, físicas, intelectuais ou sociais, pois a \textbf{brincadeira} é essencial à aprendizagem e desenvolvimento. [ ^1 ] : O termo sala de referência é preconizado na Resolução nº 5/2009, que fixa as Diretrizes Curriculares Nacionais para a Educação \textbf{Infantil}.

% matched lemmas: adulto, ajuda, areia, bater, bebê, brincadeira, brincar, brinquedo, chão, contar, conversar, creche, criança, cuidado, demonstrar, desejo, diário, gesto, imitar, infantil, interessar, limpeza, pequeno, progressivamente, pré-escola, rotina, sensorial, som, volta
\end{textsample}
